\section{二阶与三阶行列式}

\subsection{二元线性方程组与二阶行列式}

用消元法解二元线性方程组
\[
\begin{cases}
a_{11}x_1+a_{12}x_2=b_1,\\
a_{21}x_1+a_{22}x_2=b_2,
\end{cases}
\]
以$a_{22}$与$a_{12}$分别乘两方程两端再相减消去$x_2$，同理消去$x_1$，得
\[
(a_{11}a_{22}-a_{12}a_{21})x_1=b_1a_{22}-a_{12}b_2,\qquad
(a_{11}a_{22}-a_{12}a_{21})x_2=a_{11}b_2-b_1a_{21}.
\]
当$a_{11}a_{22}-a_{12}a_{21}\neq 0$时，方程组的解由四个系数与两个常数项按“两对相乘再相减”的固定格式组成。把这四个数按原位置排成二行二列的数表，并规定

\begin{definition}
由$2^2$个数排成的数表所确定的表达式
\[
\begin{vmatrix}
a_{11} & a_{12}\\
a_{21} & a_{22}
\end{vmatrix}
=a_{11}a_{22}-a_{12}a_{21}
\]
称为\textbf{二阶行列式}。数$a_{ij}$称为行列式的\textbf{元素}：第一个下标$i$为\textbf{行标}（位于第$i$行），第二个下标$j$为\textbf{列标}（位于第$j$列）；位于第$i$行第$j$列的元素称为$(i,j)$元。
\end{definition}

从$a_{11}$到$a_{22}$的实连线称为\textbf{主对角线}，从$a_{12}$到$a_{21}$的连线称为\textbf{副对角线}。于是二阶行列式等于“主对角线上两元素之积减去副对角线上两元素之积”，这一记忆法称为\textbf{对角线法则}。

记方程组的\textbf{系数行列式}$D$，把$D$的第1列换成常数项得$D_1$，把第2列换成常数项得$D_2$：
\[
D=\begin{vmatrix}a_{11}&a_{12}\\a_{21}&a_{22}\end{vmatrix},\qquad
D_1=\begin{vmatrix}b_1&a_{12}\\b_2&a_{22}\end{vmatrix},\qquad
D_2=\begin{vmatrix}a_{11}&b_1\\a_{21}&b_2\end{vmatrix},
\]
当$D\neq0$时，方程组的解可直接写出：

\keyformula{x_1=\frac{D_1}{D},\qquad x_2=\frac{D_2}{D}}

这就是克拉默法则的二阶情形（一般情形见第2章§4）。注意分子$D_j$是把$D$的\textbf{第$j$列}换成常数项列所得的行列式，列的位置不能换错。

\subsection{三阶行列式}

\begin{definition}
设有$3^2$个数排成3行3列的数表，规定
\[
\begin{vmatrix}
a_{11}&a_{12}&a_{13}\\
a_{21}&a_{22}&a_{23}\\
a_{31}&a_{32}&a_{33}
\end{vmatrix}
=a_{11}a_{22}a_{33}+a_{12}a_{23}a_{31}+a_{13}a_{21}a_{32}
-a_{13}a_{22}a_{31}-a_{11}a_{23}a_{32}-a_{12}a_{21}a_{33},
\]
称为该数表所确定的\textbf{三阶行列式}。
\end{definition}

三阶行列式含$6$项，每项都是不同行、不同列的三个元素之积再冠以正负号，规律遵循\textbf{对角线法则}：平行于主对角线的三条“实线”上三元素之积冠正号，平行于副对角线的三条“虚线”上三元素之积冠负号。

\corepoint{
对角线法则\textbf{只适用于二阶与三阶}行列式，四阶及以上不成立（四阶行列式有$4!=24$项，而“对角线”只能给出$8$项）。三阶以上必须用§3的定义或§5的展开法则计算。
}

与二元方程组类似，对三元线性方程组，若其系数行列式$D\neq0$，则未知数可用三阶行列式表示：$x_j=D_j/D$，其中$D_j$是把$D$的第$j$列换成常数项列所得的三阶行列式（见第2章§4克拉默法则）。

\section{全排列与对换}

\subsection{排列及其逆序数}

把$n$个不同的元素排成一行，称为这$n$个元素的\textbf{全排列}（简称\textbf{排列}）。所有$n$阶排列构成的集合记为$S_n$，其中元素个数即$n$阶排列的个数，记为$|S_n|$：从$n$个元素中任取一个放第一个位置有$n$种取法，再从剩下的$n-1$个中取一个放第二个位置有$n-1$种取法……最后只剩一个元素放在第$n$个位置，于是
\[
|S_n|=n(n-1)\cdots 2\cdot 1=n!.
\]

对$n$个不同元素先规定一个\textbf{标准次序}（$n$个不同自然数通常规定由小到大为标准次序）。在任一排列中，当某一对元素的先后次序与标准次序不同时，就说它构成1个\textbf{逆序}；一个排列中所有逆序的总数叫做这个排列的\textbf{逆序数}，排列$p_1p_2\cdots p_n$的逆序数记为$t(p_1p_2\cdots p_n)$。

\begin{definition}
逆序数为奇数的排列称为\textbf{奇排列}，逆序数为偶数的排列称为\textbf{偶排列}。自然排列$12\cdots n$的逆序数为$0$，是偶排列。
\end{definition}

\textbf{逆序数的计算}：不妨设元素为$1,2,\cdots,n$且由小到大为标准次序。对排列中的每个元素$a_i=p_i$，数出排在它\textbf{前面}且\textbf{比它大}的元素个数$t_i$（即$a_i$的逆序数），则
\[
t(p_1p_2\cdots p_n)=t_1+t_2+\cdots+t_n.
\]
等价地，也可数“排在它后面且比它小”的元素个数，逐元素累加；两种数法结果相同。

\subsection{对换}

在排列中，将任意两个元素对调、其余元素不动，这种作出新排列的操作叫做\textbf{对换}；将相邻两个元素对换，叫做\textbf{相邻对换}。记对换元素$a,b$为$(a\ b)$。

\begin{theorem}
一个排列中的任意两个元素对换，排列改变奇偶性。
\end{theorem}

证明思路：先证相邻对换——设排列为$\cdots ab\cdots$，对换$a,b$后，除$a,b$外其余元素的逆序数不变；当$a<b$时对换后$a$的逆序数增加1而$b$的不变，当$a>b$时$a$的不变而$b$的减少1，故排列奇偶性改变。一般对换可归结为奇数次相邻对换（先经$m$次相邻对换把$a$移到$b$左侧，再经$m+1$次把$b$移到$a$的位置，共$2m+1$次），而每次相邻对换都改变奇偶性，共改变奇数次，故奇偶性改变。

\begin{corollary}
奇排列对换成标准排列的对换次数为奇数，偶排列对换成标准排列的对换次数为偶数。
\end{corollary}

\corepoint{
本节的结论只为§3给出$n$阶行列式的一般定义、并在§4证明行列式性质时使用，属于“非必读”的理论准备。考试中的常考形式是：给定排列求逆序数、判断奇偶性；或由“对换改变奇偶性”推出“$n\ge2$时奇、偶排列各占$n!/2$个”。
}

\section{$n$阶行列式的定义}

先看三阶行列式的结构，以便推广。三阶行列式的6项中，每项都是不同行、不同列的三个元素之积，故除符号外每项可写成$a_{1p_1}a_{2p_2}a_{3p_3}$：第一个下标（行标）排成标准次序$123$，第二个下标（列标）$p_1p_2p_3$是$1,2,3$的某个排列。这样的排列共$3!=6$种，正好对应6项。再看符号：带正号的三项列标排列为$123,231,312$，都是偶排列；带负号的三项列标排列为$132,213,321$，都是奇排列。因此各项所带符号可统一写成$(-1)^{t}$，$t$为列标排列的逆序数，即
\[
\begin{vmatrix}
a_{11}&a_{12}&a_{13}\\
a_{21}&a_{22}&a_{23}\\
a_{31}&a_{32}&a_{33}
\end{vmatrix}
=\sum_{p_1p_2p_3}(-1)^{t(p_1p_2p_3)}a_{1p_1}a_{2p_2}a_{3p_3}.
\]

\begin{definition}
设有$n^2$个数排成$n$行$n$列的数表，作出表中位于不同行不同列的$n$个数的乘积，并冠以符号$(-1)^{t}$，得到形如
\[
(-1)^{t(p_1p_2\cdots p_n)}a_{1p_1}a_{2p_2}\cdots a_{np_n}
\]
的项，其中$p_1p_2\cdots p_n$为自然数$1,2,\cdots,n$的一个排列，$t(p_1p_2\cdots p_n)$为该排列的逆序数。这样的项共有$n!$个，所有这$n!$项的代数和
\[
D=\sum_{p_1p_2\cdots p_n}(-1)^{t(p_1p_2\cdots p_n)}a_{1p_1}a_{2p_2}\cdots a_{np_n}
\]
称为\textbf{$n$阶行列式}，记作
\[
D=\begin{vmatrix}
a_{11}&a_{12}&\cdots&a_{1n}\\
a_{21}&a_{22}&\cdots&a_{2n}\\
\vdots&\vdots&&\vdots\\
a_{n1}&a_{n2}&\cdots&a_{nn}
\end{vmatrix},
\]
简记作$\det(a_{ij})$，数$a_{ij}$称为行列式$D$的$(i,j)$元。
\end{definition}

按此定义作出的二阶、三阶行列式与§1用对角线法则定义的完全一致。当$n=1$时，一阶行列式$|a_{11}|=a_{11}$，注意不要与绝对值记号混淆。

由于行标也可以取遍所有排列，定义还有等价形式（行标为排列、列标为标准次序）：
$D=\sum_{q_1q_2\cdots q_n}(-1)^{t(q_1q_2\cdots q_n)}a_{q_11}a_{q_22}\cdots a_{q_nn}$。

\begin{definition}
主对角线以下（上）的元素全为$0$的行列式称为\textbf{上（下）三角形行列式}；主对角线以下和以上的元素全为$0$的行列式称为\textbf{对角行列式}。
\end{definition}

\begin{theorem}
上三角形、下三角形与对角行列式都等于其主对角线上各元素之积：
\[
\begin{vmatrix}
a_{11}&a_{12}&\cdots&a_{1n}\\
0&a_{22}&\cdots&a_{2n}\\
\vdots&\vdots&\ddots&\vdots\\
0&0&\cdots&a_{nn}
\end{vmatrix}
=\begin{vmatrix}
a_{11}&0&\cdots&0\\
a_{21}&a_{22}&\cdots&0\\
\vdots&\vdots&\ddots&\vdots\\
a_{n1}&a_{n2}&\cdots&a_{nn}
\end{vmatrix}
=a_{11}a_{22}\cdots a_{nn}.
\]
\end{theorem}

\keyformula{\begin{vmatrix}\lambda_1&&\\&\ddots&\\&&\lambda_n\end{vmatrix}=\lambda_1\lambda_2\cdots\lambda_n}

证明思路（以下三角为例）：当$i<j$时$a_{ij}=0$，故$D$中可能不为零的元素$a_{ip_i}$须满足$i\ge p_i$，又$p_1p_2\cdots p_n$是排列，只能有$p_i=i$，即可能不为零的项只有$a_{11}a_{22}\cdots a_{nn}$一项，其符号$(-1)^{t(12\cdots n)}=1$。

\corepoint{
由定义求$n$阶行列式要枚举$n!$项，只有“非零项极少”的特殊行列式（三角形、对角、副对角线型）才直接用定义。一般情形用§4的性质化三角形、或用§5的展开法则降阶。
}

\section{行列式的性质}

记$D^T$为把$D$的行换成同序数的列所得的行列式，称为$D$的\textbf{转置行列式}，即$D^T=\det(a_{ji})$。约定记号：$r_i\leftrightarrow r_j$表示对换$i,j$两行，$r_i(k)$表示第$i$行乘$k$，$r_i+kr_j$表示第$j$行的$k$倍加到第$i$行；列变换相应记作$c_i\leftrightarrow c_j$、$c_i(k)$、$c_i+kc_j$。

\begin{property}
行列式与它的转置行列式相等：$D^T=D$。
\end{property}

因此行列式的行与列地位同等：对行成立的每一条性质，对列也同样成立，反之亦然。以下性质只对行叙述。

\begin{property}
对换行列式的两行（记作$r_i\leftrightarrow r_j$），行列式变号。
\end{property}

\begin{corollary}
若行列式有两行完全相同，则此行列式等于零。
\end{corollary}

（证明：对换这两行，由性质2得$D=-D$，故$D=0$。）

\begin{property}
行列式的某一行中所有元素都乘同一数$k$，等于用数$k$乘此行列式（记作$r_i(k)$），即
\[
\begin{vmatrix}
a_{11}&a_{12}&\cdots&a_{1n}\\
\vdots&\vdots&&\vdots\\
ka_{i1}&ka_{i2}&\cdots&ka_{in}\\
\vdots&\vdots&&\vdots\\
a_{n1}&a_{n2}&\cdots&a_{nn}
\end{vmatrix}
=k\begin{vmatrix}
a_{11}&a_{12}&\cdots&a_{1n}\\
\vdots&\vdots&&\vdots\\
a_{i1}&a_{i2}&\cdots&a_{in}\\
\vdots&\vdots&&\vdots\\
a_{n1}&a_{n2}&\cdots&a_{nn}
\end{vmatrix}.
\]
\end{property}

\begin{corollary}
行列式中某一行所有元素的公因子可以提到行列式记号外面。
\end{corollary}

\begin{corollary}
行列式中若有两行元素成比例，则此行列式等于零。
\end{corollary}

\begin{property}
若行列式的某一行元素都是两数之和，例如第$i$行的元素都是两数之和，则$D$等于下列两个行列式之和：这两个行列式除第$i$行外其余各行与$D$完全相同，而第$i$行分别是这两组数。
\end{property}

\begin{property}
把行列式的某一行各元素乘同一数然后加到另一行对应的元素上去（记作$r_i+kr_j$），行列式不变：
\[
\begin{vmatrix}
\vdots&\vdots&&\vdots\\
a_{i1}+ka_{j1}&a_{i2}+ka_{j2}&\cdots&a_{in}+ka_{jn}\\
\vdots&\vdots&&\vdots\\
a_{j1}&a_{j2}&\cdots&a_{jn}\\
\vdots&\vdots&&\vdots
\end{vmatrix}
=\begin{vmatrix}
\vdots&\vdots&&\vdots\\
a_{i1}&a_{i2}&\cdots&a_{in}\\
\vdots&\vdots&&\vdots\\
a_{j1}&a_{j2}&\cdots&a_{jn}\\
\vdots&\vdots&&\vdots
\end{vmatrix}.
\]
\end{property}

性质2、3、5给出了行列式关于行的三种运算（对换、倍乘、倍加），它们与矩阵的初等行变换一一对应。由性质4还可看出：若$n$阶行列式每个元素都表示成两数之和，则可分解成$2^n$个行列式之和。

\subsection{用性质化三角形计算行列式}

利用上述运算可简化行列式的计算，常用的一条路线是\textbf{化三角形}：

\begin{enumerate}
\item 观察结构：若某一行（列）零元多、或各行（列）元素之和相等，优先对该行（列）下手；
\item 用$r_i+kr_j$把主对角线以下的元素逐个化为$0$（化上三角形），必要时先用$r_i(k)$把主元化为$1$以减少分数；
\item 化成三角形后，行列式等于主对角线上各元素之积，注意累计每次$r_i\leftrightarrow r_j$带来的负号。
\end{enumerate}

若行列式的零元分布在副对角线一侧（“副对角线型”），可先把第1行依次与第$n,n-1,\cdots,2$行对换，共$n-1$次对换，化为上三角形；再仿此处理后$n-1$行，最终共需$\frac{n(n-1)}{2}$次对换，故
\[
\begin{vmatrix}
&&&a_{1n}\\&&a_{2,n-1}&\\&\vdots&&\\a_{n1}&&&
\end{vmatrix}
=(-1)^{\frac{n(n-1)}{2}}a_{1n}a_{2,n-1}\cdots a_{n1}.
\]

\corepoint{
三处最常见的错：
（1）$r_i+kr_j$与$r_j+kr_i$结果不同，前者是“$j$行的$k$倍加到$i$行”，不能套用加法交换律写成$kr_i+r_j$；
（2）把几个运算写在一起（如$r_2-r_1,\ r_4+5r_1$）时，后一次运算作用于前一次的结果，\textbf{次序一般不能颠倒}，颠倒会得出错误结果；
（3）用性质3提出公因子时只对\textbf{一行}提，若$n$阶行列式每行都提一次公因子$k$，外面应是$k^n$。
}

\subsection{分块三角行列式与递推法}

若$A$为$m$阶方阵、$B$为$n$阶方阵，$O$为零矩阵，则
\[
\begin{vmatrix}A&O\\C&B\end{vmatrix}
=\begin{vmatrix}A&C\\O&B\end{vmatrix}
=|A|\,|B|,
\]
其中$|A|$、$|B|$分别表示相应方阵的行列式。（证明思路：对前$m$行作$r_i+\lambda r_j$把$A$化为下三角形，再对后$n$列作$c_j+\lambda c_i$把$B$化为下三角形，即得下三角形行列式，其值为$|A||B|$。）

对含参数或阶数不定的$n$阶行列式，常用\textbf{递推法}：先按某行（列）展开或对特殊结构降阶，导出$D_n$与$D_{n-1}$（或$D_{n-2}$）的递推公式，再验证$n=1$（或$n=2$）时结论成立，最后由递推式求出通项。归纳法的主要步骤——导出递推公式、检验初始情形——不可省略。

\section{行列式按行（列）展开}

一般说来，低阶行列式比高阶行列式容易计算，因此考虑用低阶行列式表示高阶行列式。

在$n$阶行列式中，把$(i,j)$元$a_{ij}$所在的第$i$行和第$j$列划去后，留下来的$n-1$阶行列式叫做$(i,j)$元$a_{ij}$的\textbf{余子式}，记作$M_{ij}$；记
\[
A_{ij}=(-1)^{i+j}M_{ij},
\]
叫做$(i,j)$元$a_{ij}$的\textbf{代数余子式}。$A_{ij}$与$M_{ij}$只差一个符号，符号由$i+j$的奇偶决定（等价地由$(i,j)$位置“正负相间”的棋盘格决定，$(1,1)$元处为正）。

\begin{lemma}
一个$n$阶行列式，如果其中第$i$行所有元素除$(i,j)$元$a_{ij}$外都为零，那么这个行列式等于$a_{ij}$与它的代数余子式的乘积，即
\[
D=a_{ij}A_{ij}.
\]
\end{lemma}

证明思路：先证$i=j=1$的情形（此时$D=a_{11}M_{11}=a_{11}A_{11}$，是分块三角行列式的特例）；一般情形把第$i$行依次与第$i-1,i-2,\cdots,1$行对换、第$j$列依次与第$j-1,\cdots,1$列对换，经$(i-1)+(j-1)$次对换把$a_{ij}$换成$(1,1)$元，符号为$(-1)^{i+j-2}=(-1)^{i+j}$，再用$i=j=1$的结果。

\begin{theorem}
行列式等于它的任一行（列）的各元素与其对应的代数余子式乘积之和，即
\[
D=a_{i1}A_{i1}+a_{i2}A_{i2}+\cdots+a_{in}A_{in}\quad(i=1,2,\cdots,n),
\]
\[
D=a_{1j}A_{1j}+a_{2j}A_{2j}+\cdots+a_{nj}A_{nj}\quad(j=1,2,\cdots,n).
\]
\end{theorem}

这个定理叫做\textbf{行列式按行（列）展开法则}。证明思路：把第$i$行元素拆成$n$个“只有一个非零元”的行列式之和（性质4），再对每个用引理。

\begin{corollary}
行列式某一行（列）的元素与另一行（列）对应元素的代数余子式乘积之和等于零，即
\[
a_{i1}A_{j1}+a_{i2}A_{j2}+\cdots+a_{in}A_{jn}=0\quad(i\neq j),
\]
\[
a_{1i}A_{1j}+a_{2i}A_{2j}+\cdots+a_{ni}A_{nj}=0\quad(i\neq j).
\]
\end{corollary}

综合定理与推论，关于代数余子式有统一的重要性质：
\keyformula{\sum_{k=1}^{n}a_{ik}A_{jk}=\begin{cases}D,&i=j,\\0,&i\neq j,\end{cases}\qquad
\sum_{k=1}^{n}a_{ki}A_{kj}=\begin{cases}D,&i=j,\\0,&i\neq j.\end{cases}}

后一式也常写成$\delta_{ij}D$（$\delta_{ij}$为Kronecker符号）。记忆口诀：“\textbf{本行乘本行余子式得$D$，异行乘异行余子式得$0$}”。这一结论是第2章用伴随矩阵求逆矩阵的基础。

\subsection{按行（列）展开的计算步骤}

\begin{enumerate}
\item 选取零元最多的一行（列）作为展开行（列）；若没有明显的零元，用$r_i+kr_j$先把该行（列）化成只有一个非零元；
\item 按该行（列）展开，把$n$阶降为$n-1$阶，注意$(-1)^{i+j}$的符号；
\item 重复上述过程，直到化为二阶或三阶行列式，用对角线法则算到底。
\end{enumerate}

若需“造零”的那一行（列）中有某个元为$1$，就以它为主元，用$r_i-r_1$等运算把同列其余元素一次消为零；若主元不为$1$，可先用$r_i(k)$把它化为$1$，避免出现分数。

\subsection{范德蒙德行列式}

称形如下式的$n$阶行列式为\textbf{范德蒙德（Vandermonde）行列式}：
\[
D_n=\begin{vmatrix}
1&1&\cdots&1\\
x_1&x_2&\cdots&x_n\\
x_1^2&x_2^2&\cdots&x_n^2\\
\vdots&\vdots&&\vdots\\
x_1^{n-1}&x_2^{n-1}&\cdots&x_n^{n-1}
\end{vmatrix}
=\prod_{1\le j<i\le n}(x_i-x_j).
\]
证明用数学归纳法：从第$n$行开始，后行减去前行的$x_1$倍，再按第1列展开并把每列的公因子$(x_i-x_1)$提出，即化为$n-1$阶范德蒙德行列式。

\corepoint{
（1）范德蒙德行列式要求“列（行）成等比”，即每列是同一个数的$0,1,\cdots,n-1$次幂；若某行缺项或幂次跳跃，不可直接套公式。
（2）公式右端是\textbf{所有}$x_i-x_j$（$i>j$）的乘积，共$\frac{n(n-1)}{2}$个因子；符号规律是大下标减小下标，写反会整体差一个$(-1)^{\frac{n(n-1)}{2}}$。
（3）$D_n=0\iff$有两个$x_i$相等，这一判据常用于“证明某行列式为零”或“求方程$D_n=0$的根”。
}
